\section{引言}

甲状腺乳头状微小癌（PTMC）总体预后较好，但中央区淋巴结转移（CentralLNM）会显著影响术式选择、淋巴结清扫策略及术后风险分层。对于术前超声提示不典型、而临床又存在潜在转移风险的患者，如何在手术前尽可能准确地完成 CentralLNM 风险评估，始终是一个兼具临床价值与技术挑战的问题。

围绕这一问题，现有研究通常从两条路径切入。第一条路径是基于术前临床变量建立风险模型，其优点是变量来源清晰、解释性较强，但在复杂影像表型面前可能存在信息不足。第二条路径是利用超声图像开展影像组学或深度学习建模，理论上能够捕获临床变量难以表征的视觉模式，但其结果容易受到样本规模、成像差异、标签噪声与融合策略的影响，导致不同研究之间结论不稳定。因此，对于 CentralLNM 这样一个术前信号有限、但临床决策价值很高的任务，更关键的问题并非“影像能否完全替代临床”，而是“影像能否在原始临床基线已经较强的前提下提供稳定且可解释的增量信息”。

本项目此前已经完成了大量实验尝试，包括原始临床回归、传统影像组学、多模态深度学习、患者级 MIL、骨干网络比较、消融验证以及三源融合探索。问题不在于结果不足，而在于现有材料更接近“实验推进过程记录”，尚未完全转换为符合临床期刊阅读习惯的论文主线。因此，本稿不再按照实验发生的先后顺序叙述，而是围绕临床读者真正关心的问题重新组织证据链。

本研究的核心目标有三点。第一，在统一患者级 5-fold 交叉验证协议下，明确原始临床基线、传统影像组学与多模态深度学习三类方法对 CentralLNM 的相对表现，并量化其相对于原始临床基线的真实增量幅度与稳定性变化。第二，解释多模态深度学习主模型为何需要患者级多实例学习（MIL）、临床主导门控融合以及临床先验约束，并通过骨干公平比较与消融验证说明其有效性。第三，进一步评估在深度模型基础上加入 ROI radiomics 后，是否能够形成稳定且公平的多源融合增益。

基于上述目标，本文将全文结果依次重组为五个层次：原始临床基线模型、传统影像组学方法、多模态深度学习融合主模型、多架构公平比较与关键模块消融，以及多源融合扩展验证。这样的组织方式更接近临床期刊的论证逻辑，也更利于区分“最终应呈现的核心证据”和“仅用于内部排查的过程性结果”。
